\examquestion{2025}{2}{Ytterbium level diagram, slowing and laser cooling}{3}
The source paper supplies a Yb level diagram, including
\(6s^2\,{}^1S_0\), \(6s6p\ P\) and \(5d6s\ D\) terms.
Distinguish strictly LS-allowed E1 transitions from weak
spin-mixed intercombination lines.
\subsection*{(a) Identify E1 arrows}
Electric dipole (E1) has odd parity and angular rank 1.
In a good LS-coupling approximation, require a parity
change, \(\Delta S=0\), and \(\Delta J=0,\pm1\)
excluding \(0\leftrightarrow0\).
For the terms visible in the supplied diagram,
\[
\begin{array}{l|l}
\text{Level pair}&\text{Strict LS E1?}\\\hline
{}^1P_1\leftrightarrow{}^1S_0&\text{yes}\\
{}^1P_1\leftrightarrow{}^1D_2&\text{yes}\\
{}^3P_1\leftrightarrow{}^3D_1&\text{yes}\\
{}^3P_1\leftrightarrow{}^3D_2&\text{yes}\\
{}^3P_1\leftrightarrow{}^3D_3&\text{no, }\Delta J=2\\
{}^3P_1\leftrightarrow{}^1S_0&
\text{spin-forbidden in pure LS coupling}
\end{array}
\]
With spin--orbit mixing the
\({}^3P_1\leftrightarrow{}^1S_0\) line becomes
a weak \emph{intercombination} E1 transition.
A \(J=0\leftrightarrow0\) line remains forbidden
for a rank-1 operator even with state mixing.
If additional \(P_J\) levels are present in the
original figure, apply the same rules to each.
See the original level diagram in the paper and
the corresponding page preview retained with
the source extractions.

\subsection*{(b) Clock transition}
A clock favours a spectrally narrow, long-lived
transition over a bright, rapidly decaying one.
Among the \(P\) terms visible in the supplied
diagram, one candidate is the much narrower
intercombination transition
\[
\boxed{6s^2\,{}^1S_0\longleftrightarrow
6s6p\,{}^3P_1\quad(\sim556~\mathrm{nm})}.
\]
It is spin-forbidden in pure LS coupling, but
spin--orbit mixing makes it weakly allowed.
This is narrower than the bright
\({}^1S_0\to{}^1P_1\) cooling line and is
a possible frequency reference. The
\emph{standard highest-precision} optical Yb
clock instead uses
\({}^1S_0\leftrightarrow{}^3P_0\) near 578 nm,
a nominally E1-forbidden \(J=0\to0\) line,
in isotopes where it can be interrogated.
Choose that state if it appears in the original
level diagram; either answer must explain
why the chosen excited state is long-lived.

\subsection*{(c) Laser-cooling transition}
Use the strong, nearly cycling
\[
\boxed{6s^2\,{}^1S_0\longleftrightarrow
6s6p\,{}^1P_1\quad(\lambda\simeq399~\mathrm{nm})}.
\]
Its large spontaneous-emission rate allows many
momentum kicks. The narrower
\({}^1S_0\to{}^3P_1\) line near 556 nm
is useful for subsequent cooling.

\subsection*{(d.i) Why the Zeeman slower needs a magnet}
As atoms slow, their Doppler shift changes,
tending to move the laser away from resonance.
A position-dependent magnetic field Zeeman-shifts
the transition and compensates the changing
Doppler shift, maintaining efficient scattering.

\subsection*{(d.ii--iii) Stopping distance and number of photons}
The data are \(M=1.73\times10^{-25}~\mathrm{kg}\),
\(T=300~\mathrm{K}\), \(\lambda=399~\mathrm{nm}\),
and excited-state lifetime \(\tau=5.5~\mathrm{ns}\).
The paper does not specify which characteristic
thermal beam speed to use. Taking the
3D rms speed as an explicit convention,
\[
v_0=\sqrt{\frac{3k_BT}{M}}=
\boxed{268.0~\mathrm{m\,s^{-1}}}.
\]
At saturation a two-level transition has
maximum photon-scattering rate \(\Gamma/2\)
where \(\Gamma=1/\tau\).
Each photon transfers mean directed momentum
\(h/\lambda\), so
\[
F_{\max}=\frac{h}{2\lambda\tau},\quad
a_{\max}=\frac{h}{2M\lambda\tau}.
\]
At \emph{half} the maximum deceleration,
\[
a=\frac{a_{\max}}2
=\frac{h}{4M\lambda\tau}
=\boxed{4.36\times10^5~\mathrm{m\,s^{-2}}}.
\]
Using \(0=v_0^2-2aL\),
\[
\boxed{L=\frac{v_0^2}{2a}=0.0823~\mathrm{m}
=8.23~\mathrm{cm}}.
\]
Momentum balance requires
\[
\boxed{N_\gamma=\frac{Mv_0}{h/\lambda}
=\frac{Mv_0\lambda}{h}\simeq2.79\times10^4}.
\]
If instead one uses the Maxwell most-probable
speed \(\sqrt{2k_BT/M}=218.8~\mathrm{m\,s^{-1}}\),
the corresponding results are
\(L=5.49~\mathrm{cm}\) and
\(N_\gamma\simeq2.28\times10^4\).
State the chosen speed definition in the exam.

\subsection*{(e.i) Red-detuned optical molasses}
An atom moving into a counter-propagating
red-detuned beam sees its light shifted toward
resonance and preferentially absorbs photons
that oppose its motion. The average force near
rest has the form \(\mathbf F\simeq-\beta\mathbf v\).

\subsection*{(e.ii) Doppler limit}
For a simple two-level transition,
\[
T_D=\frac{\hbar\Gamma}{2k_B}
=\frac{h}{4\pi k_B\tau}
=\boxed{6.94\times10^{-4}~\mathrm{K}
=0.694~\mathrm{mK}}.
\]
\solcheck{Cooling requires a bright, rapid
cycling transition; clocks require a narrow,
long-lived one. Check the factor of one-half
in the stated deceleration and the units of
acceleration, length, temperature and photon count.}
